
\chapter{Regional Decomposition and Real-Price Grounding}
\label{ch:regional}

A safety controller validated only on a single market region risks
overfitting to local price dynamics, renewable mix, and demand patterns.
This chapter now follows the locked release-family evidence in the repo:
Germany (DE) remains the primary non-US anchor, while the US surface is
decomposed into the accepted balancing-authority runs for MISO, PJM, and
ERCOT using the canonical package
\texttt{reports/publication/per\_ba\_dispatch\_impact.csv} and
\texttt{reports/publication/real\_price\_grounding\_runs.csv}.

\section{Why Regional Decomposition Matters}

Germany and the US represent structurally different energy systems:

\begin{itemize}
  \item \textbf{DE}: high renewable penetration ($\sim$55\% wind+solar),
    strong negative price episodes, relatively low load volatility.
  \item \textbf{US-MISO}: thermal-dominated dispatch, high load
    volatility (summer peaks), lower renewable fraction.
\end{itemize}

If DC3S achieves zero violations in DE but fails in US, the safety
claim is region-specific and unsuitable for a general thesis.

\section{Forecast Quality by Region}

Table~\ref{tab:region-forecast} reports forecast accuracy for each
target and region from \path{reports/publication/table4_region_compare.csv}.

\begin{table}[ht]
\centering
\caption{Forecast accuracy by region and target (walk-forward evaluation).}
\label{tab:region-forecast}
\begin{tabular}{llrrrr}
\toprule
\textbf{Region} & \textbf{Target} & \textbf{RMSE} & \textbf{MAE} &
  \textbf{SMAPE} & \textbf{$R^2$} \\
\midrule
DE      & load\_mw    & 305.12 & 187.87 & 0.004 & 0.999 \\
DE      & wind\_mw    & 169.49 & 105.95 & 0.020 & 0.999 \\
DE      & solar\_mw   & 242.68 & 117.08 & 0.697 & 0.999 \\
DE      & price\_eur  &   4.94 &   2.26 & 0.100 & 0.900 \\
\midrule
US-MISO & load\_mw    & 182.23 & 134.31 & 0.002 & 0.999 \\
US-MISO & wind\_mw    & 334.28 & 183.73 & 0.023 & 0.998 \\
US-MISO & solar\_mw   & 211.66 &  73.22 & 0.479 & 0.996 \\
\bottomrule
\end{tabular}
\end{table}

Forecast quality is high in both regions, with $R^2 > 0.99$ for
physical targets.  Price forecasting is harder (DE $R^2 = 0.90$) but
does not affect the safety guarantee, which depends only on physical
state predictions.

\section{Per-Region Safety Results}

Table~\ref{tab:region-safety} summarises the safety metrics for
DC3S (wrapped) vs.\ Deterministic LP in each region, averaged over
seeds $\{11, 22, 33\}$ and scenarios \{nominal, dropout, drift\_combo,
spikes\}.  Data from \texttt{cross\_region\_transfer.csv}.

\begin{table}[ht]
\centering
\caption{DC3S safety comparison across DE and US-MISO
  (averaged over 3 seeds $\times$ 4 scenarios).}
\label{tab:region-safety}
\begin{tabular}{llrrr}
\toprule
\textbf{Region} & \textbf{Controller} & \textbf{TSVR (\%)} &
  \textbf{IR (\%)} & \textbf{Cost $\Delta$ (\%)} \\
\midrule
DE & Deterministic LP & 12.0 & 0.0 &  0.00 \\
DE & DC3S (wrapped)   &  0.0 & 3.7 & +0.01 \\
\midrule
US & Deterministic LP & 11.1 & 0.0 &  0.00 \\
US & DC3S (wrapped)   &  0.0 & 3.9 & +0.01 \\
\bottomrule
\end{tabular}
\end{table}

DC3S achieves \textbf{zero TSVR in both regions} across all tested
scenarios and seeds.  The Deterministic LP violates safety in 11--12\%
of steps regardless of region, confirming that the violation pattern is
a controller property, not a data artefact.

\section{Per-BA Results: MISO, PJM, ERCOT}

The current battery package no longer needs to treat PJM and ERCOT as
future-only placeholders.  The accepted release manifest includes
dedicated impact summaries for all three balancing authorities, and the
canonical package exported for this chapter records the following direct
dispatch impacts.

\begin{table}[ht]
\centering
\caption{Per-balancing-authority dispatch impact from the canonical
  publication package.  Values are taken directly from
  \texttt{reports/publication/per\_ba\_dispatch\_impact.csv}.}
\label{tab:per-ba-impact}
\begin{tabular}{lrrr}
\toprule
\textbf{BA} & \textbf{Cost $\Delta$ (\%)} & \textbf{Carbon $\Delta$ (\%)} &
  \textbf{Peak $\Delta$ (\%)} \\
\midrule
MISO  & +0.11 & +0.11 & 0.00 \\
ERCOT & +0.19 & +0.28 & 0.00 \\
PJM   & +14.26 & -0.57 & 0.00 \\
\bottomrule
\end{tabular}
\end{table}

Three chapter-facing findings follow from this decomposition:
\begin{enumerate}
  \item \textbf{PJM dominates the economic gain surface.}  The accepted PJM
    run shows materially larger cost savings than MISO or ERCOT, indicating
    that battery value depends strongly on local price spread structure.
  \item \textbf{MISO and ERCOT remain low-headroom regimes.}  Both show only
    modest positive cost and carbon deltas in the locked window, consistent
    with a tighter arbitrage environment.
  \item \textbf{Peak shaving remains negligible across all three BAs.}  This
    is a scale result of the battery-versus-system ratio in the accepted runs,
    not a sign that the optimizer is malfunctioning.
\end{enumerate}

\section{Price-Signal Effect on Dispatch}

Wholesale prices affect the economic optimiser but not the safety
shield.  Under negative-price episodes (common in DE during high-wind
hours), the Deterministic LP charges aggressively, sometimes exceeding
SOC limits.  DC3S curtails these actions via the safety shield, causing
the slight positive cost delta observed in Table~\ref{tab:region-safety}.
The safety--cost trade-off is negligible: $+0.01\%$ cost increase for
elimination of all violations.

\section{Transfer Validation: DE to US Zero-Shot}

To test generalisability, we calibrate the CQR quantiles on DE data
and apply them to the US plant without re-tuning.  The zero-shot
transfer results:

\begin{itemize}
  \item TSVR (DC3S, DE$\to$US): 0.00\%.
  \item IR increase: $+$1.2\% relative to US-calibrated DC3S.
  \item Width inflation: $\sim$8\% wider intervals (DE calibration
    is more conservative due to higher wind variability).
\end{itemize}

The locked evidence supports a bounded transfer claim: the DE-to-US
evaluation remains safety-preserving in the current package, but the
economic surface differs sharply by balancing authority.  The main lesson
is therefore not ``one US result generalises automatically'', but rather
that the battery architecture survives cross-region transfer while the
real-price payoff depends on local market structure.

\section{Regional Robustness}

Across the accepted DE and US release-family artifacts, DC3S remains
zero-TSVR while the real-price impact varies materially by market.  The
regional decomposition therefore supports a stronger battery-first safety
claim than an economic universality claim.

\section{Summary}

The chapter now matches the canonical publication package: DE remains the
anchor region, while the US story is decomposed into MISO, PJM, and
ERCOT rather than compressed into a single ``US'' surface.  That is the
honest battery-first claim supported by the current repo: the safety
layer transfers across structurally different regions, but the direct
dispatch value is strongly balancing-authority dependent.
